% !TeX root = ../../Book1.tex
% 纲要标题：### 20.3 从表示到齐性向量丛
% 纲要位置：几何学计划纲要.md，第 905 行。
% 正文状态：初稿；尚未正式审核。

% 纲要内容开始
% 
% * 稳定子表示与关联丛
% * 截面的等变函数模型
% * 截面表示与 Frobenius 互反
% * 复射影直线上的齐性线丛
% 
% 纲要内容结束

\section{从表示到齐性向量丛}
\label{g1:sec:2003}

一个群的表示在固定向量空间中改变向量；一个齐性空间上的向量丛则让向量空间随底点移动。两者的联系来自稳定子：先规定稳定子怎样作用于一个基点上的纤维，再用全群把这条纤维搬到各处。所得截面可以重新写成群上的函数，群表示又表现为函数的平移。本节使用\cref{g1:exm:homogeneous-principal-bundle,g1:thm:associated-vector-bundle}中的右主丛与关联丛约定，重点说明齐性空间额外提供的左作用以及它怎样控制截面。

\subsection{稳定子表示与关联丛}
\label{g1:subsec:2003:01}
设 \(G\) 为有限维实 Lie 群，\(H\subseteq G\) 为闭子群，\(o=eH\in G/H\)。左作用下，\(o\) 的稳定子就是 \(H\)。若向量丛 \(E\to G/H\) 带有覆盖左作用的光滑 \(G\) 作用，并且每个群元素在纤维之间作线性映射，那么 \(H\) 保持 \(E_o\)，从而给出一个纤维上的表示。这种带作用的向量丛称为\term[qixingxiangliangcong]{齐性向量丛}[homogeneous vector bundle]。下面的构造说明，这一个稳定子表示已经决定整个丛及其 \(G\) 作用。

\begin{proposition}[稳定子表示构造齐性向量丛]\label{g1:prop:homogeneous-associated-bundle}
	设 \(G\) 为有限维实 Lie 群，\(H\) 为闭子群，\(W\) 为有限维复向量空间，\(\rho:H\to\operatorname{GL}(W)\) 为光滑表示。则
	\[
		E_\rho=G\times_H W\longrightarrow G/H,\qquad
		[g,w]\longmapsto gH,
		\qquad [gh,w]=[g,\rho(h)w]
	\]
	是光滑复向量丛，且 \(a\cdot[g,w]=[ag,w]\) 定义了齐性向量丛结构。反之，每个有限秩光滑复齐性向量丛都与其基点纤维 \(W=E_o\) 的稳定子表示所构造的 \(G\times_H W\) 自然地 \(G\) 等变同构。实向量丛有相同结论。
\end{proposition}
\begin{proof}
	\(H\) 闭保证 \(G/H\) 是 Hausdorff 光滑流形，而且 \(G\to G/H\) 有光滑局部截面，是右主 \(H\) 丛。将\cref{g1:thm:associated-vector-bundle}用于这个主丛，得到所列关联丛；其等价关系来自右作用 \((g,w)\cdot h=(gh,\rho(h)^{-1}w)\)。左乘与这个右作用交换，因为 \(a(gh)=(ag)h\)，所以 \([g,w]\mapsto[ag,w]\) 良定义。在商映射之前它是光滑映射，局部平凡化又说明下降后的作用光滑，并且逐纤维复线性。

	反过来，设已有这样的 \(E\)，令 \(\rho(h)w=h\cdot w\)、\(w\in E_o\)。作用的光滑性及纤维线性性使 \(\rho\) 成为光滑线性表示。定义
	\[
		\Psi:G\times_H E_o\longrightarrow E,\qquad
		\Psi([g,w])=g\cdot w.
	\]
	它不随代表改变，因为 \((gh)\cdot w=g\cdot\rho(h)w\)。在 \(gH\) 上，群元素 \(g\) 将 \(E_o\) 线性同构地送到 \(E_{gH}\)，所以 \(\Psi\) 逐纤维双射。若 \(s:U\to G\) 是主丛局部截面，则 \(\Psi\) 的局部表达为 \((x,w)\mapsto s(x)\cdot w\)，其逆为 \(v\mapsto(p(v),s(p(v))^{-1}\cdot v)\)。这两式都光滑，故 \(\Psi\) 为丛同构。最后 \(\Psi(a\cdot[g,w])=ag\cdot w=a\cdot\Psi([g,w])\)，证明等变性。
\end{proof}

这里分类的是带指定 \(G\) 作用的丛。忘掉作用以后，不同的稳定子表示可能给出同构的普通向量丛。若 \(A:W\to W'\) 满足 \(A\rho(h)=\rho'(h)A\)，则 \([g,w]\mapsto[g,Aw]\) 定义一个等变丛映射；反过来，任何覆盖 \(G/H\) 恒等映射的等变丛映射都由其在 \(o\) 处的限制确定。这是因为它必须把 \(g\cdot w\) 送到 \(g\cdot A(w)\)。

\ProseParagraphBreak
一般表示给出的构造仍称关联丛；只有结构群的伴随表示给出的特例称为伴随丛。对于当前主 \(H\) 丛 \(G\to G/H\)，\cref{g1:def:adjoint-bundle}中的记号为
\[
	\operatorname{ad}(G\to G/H)=G\times_{\operatorname{Ad}|_H}\mathfrak h.
\]
这条丛的纤维是 \(\mathfrak h\)，不是 \(\mathfrak g\)。表示 \(\rho\) 决定了纤维中的对象及其转换规律，不能把一般关联构造中的表示都换成 \(\operatorname{Ad}\)。

\begin{example}\label{g1:exm:homogeneous-tangent-associated}
	设 \(G\) 为有限维实 Lie 群，\(H\) 为闭子群。\cref{g1:prop:homogeneous-tangent-isotropy}将 \(\mathrm{T}_o(G/H)\) 识别为 \(\mathfrak g/\mathfrak h\)，稳定子作用为 \(X+\mathfrak h\mapsto\operatorname{Ad}_hX+\mathfrak h\)。因此齐性切丛为
	\[
		\mathrm{T}(G/H)\cong G\times_H(\mathfrak g/\mathfrak h),\qquad
		[g,X+\mathfrak h]\longmapsto
		\left.\frac{\mathrm{d}}{\mathrm{d}t}\right|_{t=0}g\exp(tX)H.
	\]
	在基点，商空间同构由商映射的微分给出；在其他点用左平移搬运，正是前一命题的构造。若 \(G\) 紧，取伴随不变内积后，\(\mathfrak m=\mathfrak h^\perp\) 在 \(\operatorname{Ad}(H)\) 下不变，且 \(\mathfrak m\to\mathfrak g/\mathfrak h\) 是 \(H\) 等变线性同构，故也可用 \(G\times_H\mathfrak m\) 表示切丛。稳定子表示作用于商空间或正交补；它一般不同于作用在 \(\mathfrak h\) 上的伴随表示。
\end{example}

\subsection{截面的等变函数模型}
\label{g1:subsec:2003:02}
要在 \(gH\) 上指定截面值，可以先写成 \([g,f(g)]\)。同一个底点也能写成 \(ghH\)，所以函数值不能任意变化；它必须恰好补偿代表 \(g\) 的改变。这恢复了第十八章的等变函数关系，并且让截面变成可以直接平移、求导和积分的对象。

\begin{proposition}[齐性丛截面的函数模型]\label{g1:prop:equivariant-functions-sections}
	设 \(G\) 为有限维实 Lie 群，\(H\) 为闭子群，\(\rho:H\to\operatorname{GL}(W)\) 为有限维复光滑表示。令 \(E_\rho=G\times_H W\to G/H\)，采用 \([gh,w]=[g,\rho(h)w]\) 的商约定，并令
	\[
		C^\infty(G,W)_\rho
		=\{f\in C^\infty(G,W):f(gh)=\rho(h)^{-1}f(g)
		\;\text{对所有 }g\in G,h\in H\}.
	\]
	则线性映射
	\[
		f\longmapsto\sigma_f,\qquad
		\sigma_f(gH)=[g,f(g)]
	\]
	给出 \(C^\infty(G,W)_\rho\cong\Gamma^\infty(G/H,E_\rho)\)。若主丛局部截面满足 \(s_j=s_i h_{ij}\)，则相应局部分量 \(f_i=f\circ s_i\) 满足 \(f_i=\rho(h_{ij})f_j\)。
\end{proposition}
\begin{proof}
	这属于\cref{g1:prop:associated-sections}的情形：主丛为 \(G\to G/H\)，其右作用是 \(g\mapsto gh\)。具体地，
	\[
		[gh,f(gh)]=[g,\rho(h)f(gh)]=[g,f(g)],
	\]
	所以所给公式不依赖余类代表。在局部截面 \(s_i\) 下，\(\sigma_f(x)=[s_i(x),f(s_i(x))]\)，故截面光滑。

	反过来，对每个 \(g\)，映射 \(w\mapsto[g,w]\) 是到纤维 \((E_\rho)_{gH}\) 的线性同构，所以 \(\sigma(gH)=[g,f(g)]\) 唯一确定 \(f(g)\)。比较 \(g\) 与 \(gh\) 得到所述等变性；在 \(g=s_i(x)h\) 的局部坐标中，有 \(f(g)=\rho(h)^{-1}f_i(x)\)，所以它光滑。两种构造逐点互逆。最后，\(f_j(x)=f(s_i(x)h_{ij}(x))=\rho(h_{ij}(x))^{-1}f_i(x)\)，即所写的转换式。
\end{proof}

若 \(\rho\) 是平凡的一维表示，等变条件退化为 \(f(gh)=f(g)\)，也就是普通函数从 \(G\) 下降到 \(G/H\) 的条件。非平凡表示允许函数沿同一条纤维变化，但变化已经由 \(\rho\) 完全规定。\cref{g1:fig:associated-section-descent}把这个条件写成两个商映射之间的相容关系：函数与群元素配在一起后，才能下降为截面。

% !TeX root = ../Book1.tex
\begin{figure}[htbp]
	\centering
	\begin{tikzcd}[row sep=large,column sep=huge]
		G \arrow[r,"{\widetilde\sigma:g\mapsto(g,f(g))}"] \arrow[d,"q"']
		& G\times W \arrow[d,"Q"] \\
		G/H \arrow[r,"{\sigma:gH\mapsto[g,f(g)]}"']
		& E_\rho
	\end{tikzcd}
	\[
		\widetilde\sigma(gh)
		=(gh,f(gh))
		=(gh,\rho(h)^{-1}f(g))
		=\widetilde\sigma(g)\cdot h.
	\]
	\caption[等变函数下降为关联丛截面]{上方映射将每个群元素及其函数值配在一起。右等变条件保证它把同一条 \(H\) 轨道送到同一条 \(H\) 轨道，因而下降为下方的截面。两个竖箭头分别是主丛商映射及关联丛商映射。}
	\label{g1:fig:associated-section-descent}
\end{figure}


\subsection{截面表示与 Frobenius 互反}
\label{g1:subsec:2003:03}
群元素 \(a\) 把 \(x\) 处的纤维送到 \(ax\) 处。要得到一条新截面在 \(x\) 处的值，应先去 \(a^{-1}x\) 处读取原截面，再由 \(a\) 将这个值送回来。这给出公式
\[
	(U_a\sigma)(x)=a\cdot\sigma(a^{-1}x).
\]
逆元出现于底点，正是为了使结果仍落在 \(E_x\)，同时使 \(U_aU_b=U_{ab}\)。

\begin{proposition}[截面上的左表示]\label{g1:prop:section-left-representation}
	设 \(G\) 为有限维实 Lie 群，\(H\) 为闭子群，\(\rho:H\to\operatorname{GL}(W)\) 为有限维复光滑表示。令 \(E_\rho=G\times_H W\)，采用 \([gh,w]=[g,\rho(h)w]\) 的商约定。公式 \((U_a\sigma)(x)=a\cdot\sigma(a^{-1}x)\) 在 \(\Gamma^\infty(G/H,E_\rho)\) 上给出一个线性表示，其中总空间的作用为 \(a\cdot[g,w]=[ag,w]\)。以 \(\sigma_f(gH)=[g,f(g)]\) 识别满足 \(f(gh)=\rho(h)^{-1}f(g)\) 的光滑函数，即\cref{g1:prop:equivariant-functions-sections}的函数模型。在此模型中，它是
	\begin{equation}\label{g1:eq:homogeneous-section-left-action}
		(U_af)(g)=f(a^{-1}g).
	\end{equation}
	若 \(G\) 紧，以 \(\mu_G\) 为 \(G\) 上满足 \(\mu_G(G)=1\) 的左右不变 Haar 测度，并在 \(W\) 上选取 \(H\) 不变 Hermitian 内积，则
	\[
		\langle\sigma_f,\sigma_k\rangle_{L^2}
		=\int_G\langle f(g),k(g)\rangle_W\,\mathrm{d}\mu_G(g)
	\]
	给出截面空间上的正定内积，各 \(U_a\) 保持此内积，因而延拓为其 Hilbert 完备化上的酉算子。
\end{proposition}
\begin{proof}
	截面公式的纤维位置已核对；光滑性来自群作用及截面的光滑性。直接计算
	\[
		(U_aU_b\sigma)(x)
		=a\cdot\bigl(b\cdot\sigma(b^{-1}a^{-1}x)\bigr)
		=(ab)\cdot\sigma((ab)^{-1}x).
	\]
	单位元作用为恒等，所以确为表示。若 \(x=gH\)，则
	\[
		(U_a\sigma_f)(gH)
		=a\cdot[a^{-1}g,f(a^{-1}g)]
		=[g,f(a^{-1}g)],
	\]
	这就得到函数模型。左平移保持右等变条件，因为
	\(f(a^{-1}gh)=\rho(h)^{-1}f(a^{-1}g)\)。

	紧性使 \(H\) 也紧，所以可以平均一个 Hermitian 内积使其 \(H\) 不变。于是
	\[
		\langle f(gh),k(gh)\rangle_W=\langle f(g),k(g)\rangle_W,
	\]
	被积函数下降为底空间上的函数；积分也等于对 \(G/H\) 上由 \(\mu_G\) 推前的概率测度积分。若 \(f\ne0\)，则它在某个开集上范数有正下界，而归一化光滑体积在非空开集上为正，故所得内积正定。左不变性及代换 \(g=au\) 给出
	\[
		\int_G\langle f(a^{-1}g),k(a^{-1}g)\rangle_W\,\mathrm{d}\mu_G(g)
		=\int_G\langle f(u),k(u)\rangle_W\,\mathrm{d}\mu_G(u).
	\]
	因此每个 \(U_a\) 是等距线性算子，唯一延拓到完备化，延拓后的逆为 \(U_{a^{-1}}\)，故为酉算子。
\end{proof}

截面空间通常是无限维的。若想在其中寻找一个指定有限维表示 \(V\)，直接构造每一条截面会显得繁重；等变性允许只观察基点 \(o\) 上的值。以下互反关系将问题化为两个有限维稳定子表示之间的线性映射。

\begin{theorem}[光滑截面的 Frobenius 互反]\label{g1:thm:smooth-frobenius-reciprocity}
	设 \(G\) 为有限维实 Lie 群，\(H\) 为闭子群，\(\pi:G\to\operatorname{GL}(V)\)、\(\rho:H\to\operatorname{GL}(W)\) 为有限维复光滑表示。令 \(E_\rho=G\times_H W\)，采用 \([gh,w]=[g,\rho(h)w]\) 的商约定；对光滑截面 \(\sigma\)，以 \(\sigma(gH)=[g,f_\sigma(g)]\) 定义其等变函数。截面上取 \(G\) 作用 \((U_a\sigma)(x)=a\cdot\sigma(a^{-1}x)\)，其中 \(a\cdot[g,w]=[ag,w]\)。则求值于单位元给出自然线性同构
	\begin{equation}\label{g1:eq:smooth-frobenius-reciprocity}
		\operatorname{Hom}_G\bigl(V,\Gamma^\infty(G/H,E_\rho)\bigr)
		\cong\operatorname{Hom}_H(V|_H,W).
	\end{equation}
	这里两侧的 \(\operatorname{Hom}\) 都指复线性的交织映射。若左侧映射为 \(\Phi\)，则对应的 \(A\) 是
	\(A(v)=f_{\Phi(v)}(e)\)；反过来，\(A\) 对应的截面由
	\begin{equation}\label{g1:eq:frobenius-section-formula}
		f_v(g)=A\bigl(\pi(g^{-1})v\bigr),\qquad
		\Phi_A(v)(gH)=[g,f_v(g)]
	\end{equation}
	给出。
\end{theorem}
\begin{proof}
	设 \(\Phi\) 为 \(G\) 交织映射。对于 \(h\in H\)，先用 \(\Phi\) 的交织性，再用函数模型及右等变性，得到
	\[
		A(\pi(h)v)
		=f_{U_h\Phi(v)}(e)
		=f_{\Phi(v)}(h^{-1})
		=\rho(h)f_{\Phi(v)}(e)=\rho(h)A(v).
	\]
	所以 \(A\) 确为 \(H\) 交织映射。

	反过来，给定 \(A\in\operatorname{Hom}_H(V|_H,W)\)，\eqref{g1:eq:frobenius-section-formula}中的函数光滑，且
	\[
		f_v(gh)=A\bigl(\pi(h^{-1})\pi(g^{-1})v\bigr)
		=\rho(h)^{-1}f_v(g).
	\]
	所以它给出光滑截面，且线性依赖于 \(v\)。对 \(a\in G\)，
	\[
		f_{\pi(a)v}(g)=A\bigl(\pi(g^{-1}a)v\bigr)
		=f_v(a^{-1}g)=(U_af_v)(g),
	\]
	因此 \(\Phi_A\) 为 \(G\) 交织映射。其在单位元的值为 \(A(v)\)，证明一个方向的复合为恒等。若从 \(\Phi\) 出发再构造，则对任意 \(g\)，
	\[
		f_{\Phi(v)}(g)
		=f_{U_{g^{-1}}\Phi(v)}(e)
		=f_{\Phi(\pi(g^{-1})v)}(e)
		=A\bigl(\pi(g^{-1})v\bigr).
	\]
	故重建的截面就是原截面，另一方向的复合也为恒等。
\end{proof}

例如，取 \(V=\mathbb C\) 为平凡表示，得到
\[
	\Gamma^\infty(G/H,E_\rho)^G\cong W^H.
\]
群不变截面对应函数模型中的常函数，而这样的常函数必须取值于稳定子的不动子空间。若 \(V\) 不可约，任一非零交织映射 \(V\to\Gamma^\infty(E_\rho)\) 的核都是不变子空间，故它单射；于是 \(\operatorname{Hom}_H(V|_H,W)\ne0\) 正是这个不可约表示能够出现为截面子空间的判据。它不意味着全部光滑截面都是有限多个不可约分量的和；无限展开仍须另行说明收敛方式。

\subsection{复射影直线上的齐性线丛}
\label{g1:subsec:2003:04}
现在取 \(G=\operatorname{SU}(2)\) 及极大环面
\[
	T=\{t_\theta=\operatorname{diag}(e^{\mathrm{i}\theta},e^{-\mathrm{i}\theta}):\theta\in\mathbb R\}.
\]
对整数 \(m\)，令 \(\chi_m(t_\theta)=e^{\mathrm{i}m\theta}\)，并记
\[
	E_m=\operatorname{SU}(2)\times_{\chi_m}\mathbb C.
\]
这些线丛把上一章环面的整数权转变为球面上的具体几何对象。\symidx{\(E_m\)：球面上权为 \(m\) 的关联线丛}我们将直接核对符号：在本书的右主丛约定下，权一对应射影直线自身组成的线丛，而其对偶对应权负一。

\begin{proposition}[球面上的齐性线丛]\label{g1:prop:su2-hopf-associated-lines}
	设 \(T=\{t(\theta)=\operatorname{diag}(e^{\mathrm{i}\theta},e^{-\mathrm{i}\theta}):\theta\in\mathbb R\}\subseteq\operatorname{SU}(2)\)，\(m\in\mathbb Z\)，\(\chi_m(t(\theta))=e^{\mathrm{i}m\theta}\)，\(E_m=\operatorname{SU}(2)\times_{\chi_m}\mathbb C\)，采用 \([gt,z]=[g,\chi_m(t)z]\) 的商约定。由 \(gT\mapsto\mathbb C\,g(1,0)^{\mathsf T}\) 得到自然微分同胚 \(\operatorname{SU}(2)/T\cong\mathbb{CP}^1\cong S^2\)。在这个识别下，\(E_1\) 与线丛
	\[
		L=\{(\ell,v):\ell\in\mathbb{CP}^1,\ v\in\ell\}\longrightarrow\mathbb{CP}^1
	\]
	等变同构。记 \(L=\mathcal O(-1)\)，以对偶及张量幂定义 \(\mathcal O(k)\)、\(k\in\mathbb Z\)，则有光滑等变同构 \(E_m\cong\mathcal O(-m)\)。在两张射影图
	\[
		U_0:\ z=z_1/z_0,\qquad U_1:\ w=z_0/z_1
	\]
	上，可取局部标架 \(\epsilon_0^{(m)},\epsilon_1^{(m)}\) 使交叠处
	\begin{equation}\label{g1:eq:projective-line-associated-transition}
		w=1/z,\qquad
		\epsilon_1^{(m)}=w^m\epsilon_0^{(m)},\qquad
		a_0(z)=z^{-m}a_1(1/z)
	\end{equation}
	成立，其中最后一式是同一截面的局部分量转换。
\end{proposition}
\begin{proof}
	令 \(e_0=(1,0)^{\mathsf T}\)。映射 \(g\mapsto g\mathbb Ce_0\) 将 \(\operatorname{SU}(2)\) 送到 \(\mathbb{CP}^1\)。对任意单位向量 \(u=(u_0,u_1)^{\mathsf T}\)，矩阵
	\[
		g_u=\begin{pmatrix}u_0&-\overline{u_1}\\u_1&\overline{u_0}\end{pmatrix}
	\]
	属于 \(\operatorname{SU}(2)\)，并且 \(g_ue_0=u\)，所以作用传递。保持 \(\mathbb Ce_0\) 的特殊酉矩阵必须对角，恰为 \(T\)。由\cref{g1:prop:transitive-action-model}，轨道对应给出所述光滑商模型。

	取第一列还把 \(\operatorname{SU}(2)\) 微分同胚地识别为 \(S^3\)，逆映射就是 \(u\mapsto g_u\)；右乘 \(t_\theta\) 使第一列乘以 \(e^{\mathrm{i}\theta}\)。因此当前主丛正是\cref{g1:exm:hopf-principal-curvature}的 Hopf 丛。该例中的球极坐标同胚在这里写成
	\[
		[1:z]\longmapsto
		\frac{(2\operatorname{Re}z,2\operatorname{Im}z,1-|z|^2)}{1+|z|^2}.
	\]
	它在另一图 \([w:1]\)、\(w=1/z\) 上光滑延拓到南极，逆映射也由两张球极坐标给出，故 \(\mathbb{CP}^1\cong S^2\)。

	定义
	\[
		E_1\longrightarrow L,\qquad
		[g,c]\longmapsto(g\mathbb Ce_0,cge_0).
	\]
	若把 \([g,c]\) 换成 \([gt_\theta,c]\)，输出的向量变成 \(ce^{\mathrm{i}\theta}ge_0\)，与 \([g,e^{\mathrm{i}\theta}c]\) 的输出相同，所以公式良定义。它逐纤维是线性同构，在局部截面中及其逆都光滑，并与 \(\operatorname{SU}(2)\) 左作用相容。因此 \(E_1\cong L\)。对于任意整数 \(m,n\)，使用同一个 \(g\) 表示同一纤维中的元素，可定义
	\[
		E_m\otimes E_n\longrightarrow E_{m+n},\qquad
		[g,c]\otimes[g,d]\longmapsto[g,cd].
	\]
	改变代表 \(g\mapsto gt_\theta\) 时，两因子的分量分别乘以 \(e^{-\mathrm{i}m\theta}\)、\(e^{-\mathrm{i}n\theta}\)，乘积正好按权 \(m+n\) 变化，所以该映射良定义且逐纤维可逆。又有 \(E_0\cong\mathbb{CP}^1\times\mathbb C\)，于是 \(E_{-m}\cong E_m^*\)。从 \(m=1\) 出发便得全部 \(E_m\cong\mathcal O(-m)\)。

	最后计算局部转换。两张图上的主丛截面可取
	\[
		s_0(z)=\frac1{\sqrt{1+|z|^2}}
		\begin{pmatrix}1&-\overline z\\z&1\end{pmatrix},\qquad
		s_1(w)=\frac1{\sqrt{1+|w|^2}}
		\begin{pmatrix}w&-1\\1&\overline w\end{pmatrix}.
	\]
	它们的第一列分别代表 \([1:z]\)、\([w:1]\)。在交叠处，直接比较第一列得到
	\[
		s_1(1/z)=s_0(z)
		\operatorname{diag}\!\left(\frac{|z|}{z},\frac{z}{|z|}\right).
	\]
	第二列由第一列及特殊酉条件唯一确定，所以这也是完整矩阵的等式。故 \(E_m\) 的单位局部标架 \(b_i^{(m)}=[s_i,1]\) 满足
	\(b_1^{(m)}=(|z|/z)^m b_0^{(m)}\)。为了得到更简单的有理转换，取
	\[
		\epsilon_0^{(m)}=(1+|z|^2)^{m/2}b_0^{(m)},\qquad
		\epsilon_1^{(m)}=(1+|w|^2)^{m/2}b_1^{(m)}.
	\]
	由于 \(1+|1/z|^2=(1+|z|^2)/|z|^2\)，得到 \(\epsilon_1^{(m)}=z^{-m}\epsilon_0^{(m)}=w^m\epsilon_0^{(m)}\)。对于 \(m=1\)，它们在线丛 \(L\) 中正是 \(e_0(z)=(1,z)\)、\(e_1(w)=(w,1)\)。同一截面写成 \(\epsilon_0^{(m)}a_0=\epsilon_1^{(m)}a_1\)，便得最后的分量关系。
\end{proof}

局部标架的转换与截面分量的转换方向相反。比如 \(m=-n\)、\(n\ge0\) 时，\(E_{-n}=\mathcal O(n)\) 的分量必须满足 \(a_0(z)=z^n a_1(1/z)\)。这个条件将使下一节的二元齐次多项式恰好粘成全局截面；对于 \(E_n=\mathcal O(-n)\)，相同的多项式规则则不成立。

\begin{proposition}[这些线丛的非平凡性]\label{g1:prop:projective-line-bundle-winding}
	设 \(m,n\in\mathbb Z\)，\(T=\{t(\theta)=\operatorname{diag}(e^{\mathrm{i}\theta},e^{-\mathrm{i}\theta}):\theta\in\mathbb R\}\subseteq\operatorname{SU}(2)\)。对 \(j=m,n\)，令 \(\chi_j(t(\theta))=e^{\mathrm{i}j\theta}\)，\(E_j=\operatorname{SU}(2)\times_{\chi_j}\mathbb C\to\operatorname{SU}(2)/T\cong\mathbb{CP}^1\)，采用 \([gt,z]=[g,\chi_j(t)z]\) 的商约定。则在 \(\mathbb{CP}^1\) 上，\(E_m\) 作为光滑复线丛平凡，当且仅当 \(m=0\)。更一般地，\(E_m\) 与 \(E_n\) 作为覆盖底空间恒等映射的光滑复线丛同构，当且仅当 \(m=n\)。
\end{proposition}
\begin{proof}
	\(m=0\) 的平凡化已经构造。若 \(E_m\) 平凡，就有处处非零的光滑截面，其在 \(U_0,U_1\) 的分量 \(a_0,a_1\) 处处非零。在赤道 \(z=e^{\mathrm{i}\theta}\)、\(w=e^{-\mathrm{i}\theta}\) 上，转换关系给出
	\[
		a_0(e^{\mathrm{i}\theta})
		=e^{-\mathrm{i}m\theta}a_1(e^{-\mathrm{i}\theta}).
	\]
	任意非零复数环路先除以模长，就成为圆周环路；\cref{g1:thm:circle-fundamental}用连续角度的总增量定义其绕数。环路相乘时角度相加，所以绕数相加；反向参数化使绕数变号。若环路能在圆盘上保持非零地延拓，则径向收缩给出零同伦，绕数为零。

	这里 \(a_0\) 与 \(a_1\) 分别在 \(|z|\le1\)、\(|w|\le1\) 的圆盘上非零延拓，所以上式左边的绕数为零，右边第二因子的绕数也为零。第一因子的绕数为 \(-m\)，故 \(0=-m\)，只能有 \(m=0\)。最后，若 \(E_m\cong E_n\)，则 \(E_m^*\otimes E_n\cong E_{n-m}\) 有由该同构给出的处处非零截面，因而平凡；刚才的结论迫使 \(m=n\)。反向取恒等同构即可。
\end{proof}

这些线丛即使非平凡，也有大量光滑截面：在一张图内取紧支集光滑函数乘以局部标架，再在图外延拓为零即可。非平凡性排除的是处处非零的截面。另一方面，\(E_m\) 在基点的稳定子表示为 \(\chi_m\)，所以 \(m\ne0\) 时没有非零的 \(\operatorname{SU}(2)\) 不变截面。普通光滑截面的存在、不变截面的存在和全局标架的存在，是三个不同的问题。
% 纲要条目排为编号小节标题，并进入目录与书签；说明段仍保留于纲要注释中。
